\section{Representação de uma solução}

Uma solução $S=(R_1,\ldots,R_k)$ é uma lista de rotas, em que cada posição $j$ da lista corresponde a um caminhão (veículo) e $R_j$ é a rota que ele percorre, uma lista ordenada de clientes, $R=(v_1,\ldots,v_m)$. O depósito (nó $0$) não é armazenado: ele é implícito no início e no fim, de modo que a rota percorrida é $0\rightarrow v_1\rightarrow\cdots\rightarrow v_m\rightarrow 0$. Rotas vazias são descartadas, e o número de rotas não pode exceder o número de veículos da instância, obtido do sufixo \texttt{-k} do nome do arquivo.

Uma solução é \emph{viável} quando atende todos os clientes exatamente uma vez, usa no máximo $k$ rotas e respeita a capacidade em cada rota. Essa verificação é independente do custo: uma solução de custo baixo pode ser inválida se omitir ou repetir um cliente, ou exceder a capacidade. Por isso toda solução é validada antes de ser aceita como melhor solução.

\textbf{Exemplo.} Solução obtida pelo GRASP clássico ($\alpha=0.25$, 100 iterações, semente 42) na instância \texttt{A-n32-k5} (31 clientes, 5 veículos, capacidade 100; no arquivo, o depósito é o nó 1). Seu custo é 808, com gap de 3.06\% em relação ao ótimo conhecido (784):
{\small
\begin{align*}
S=\big(&[25,28],\ [21,6,30,19,12,5,24,4,3,7],\ [15,29,9,10,23,16,11,26],\\
&[31,17,2,13],\ [22,32,20,18,14,8,27]\big).
\end{align*}}
Cada posição da lista $S$ corresponde a um caminhão: o primeiro caminhão atende os clientes $[25,28]$, o segundo atende $[21,6,30,\ldots,7]$, e assim por diante, totalizando os 5 veículos da instância. Dentro de cada rota, a ordem dos clientes é a ordem de visita. As cargas dos caminhões são 44, 98, 98, 72 e 98, todas dentro da capacidade, e cada um dos 31 clientes aparece uma única vez. A primeira rota, por exemplo, percorre $\text{depósito}\rightarrow 25\rightarrow 28\rightarrow\text{depósito}$.

O custo de uma rota é
\begin{equation}
C(R)=d_{0v_1}+\sum_{i=1}^{m-1}d_{v_iv_{i+1}}+d_{v_m0},
\end{equation}
e a função objetivo, a ser minimizada, é $F(S)=\sum_{j=1}^{k}C(R_j)$. O gap em relação ao ótimo conhecido $F^*$ é
\begin{equation}
\text{gap}(S)=100\left(\frac{F(S)}{F^*}-1\right).
\end{equation}
A matriz de distâncias é calculada uma vez por instância e reutilizada na construção, na validação e nos movimentos da busca local.
